\documentclass[UTF8,fontset=none,11pt,a4paper]{ctexart}
\usepackage[left=21mm,right=21mm,top=27mm,bottom=22mm,headheight=26pt,headsep=6mm]{geometry}
\usepackage{fontspec}
\setmainfont{texgyretermes-regular.otf}[BoldFont=texgyretermes-bold.otf,ItalicFont=texgyretermes-italic.otf]
\setCJKmainfont{Songti.ttc}[Path=/System/Library/Fonts/Supplemental/,FontIndex=6,AutoFakeBold=2]
\setCJKsansfont{Songti.ttc}[Path=/System/Library/Fonts/Supplemental/,FontIndex=6,AutoFakeBold=2]
\usepackage[table]{xcolor}
\definecolor{TableHead}{HTML}{EEF0ED}
\definecolor{Muted}{HTML}{666C68}
\usepackage{booktabs,longtable,array,fancyhdr,needspace,xurl}
\usepackage[unicode,hidelinks]{hyperref}
\urlstyle{same}
\linespread{1.32}
\setlength{\parindent}{2em}
\setlength{\parskip}{3pt}
\setlength{\LTpre}{8pt}\setlength{\LTpost}{10pt}
\ctexset{section={format=\fontsize{15}{20}\selectfont\bfseries,beforeskip=18pt,afterskip=8pt},subsection={format=\fontsize{12}{16}\selectfont\bfseries,beforeskip=12pt,afterskip=6pt},subsubsection={format=\normalsize\bfseries,beforeskip=9pt,afterskip=4pt}}
\pagestyle{fancy}\fancyhf{}
\renewcommand{\headrulewidth}{.3pt}
\renewcommand{\footrulewidth}{0pt}
\fancyhead[L]{\parbox[b]{\linewidth}{\fontsize{8}{10}\selectfont\color{Muted}汽车贷款保险：已有保单与新增保障记录·本次账户资料及办理记录}}
\fancyfoot[C]{\small\color{Muted}\thepage}
\hypersetup{pdftitle={本次账户资料及办理记录},pdfsubject={仅供模拟使用·汽车贷款保险：已有保单与新增保障记录·13-03},pdfauthor={金融决策模拟材料}}
\begin{document}
{\noindent\fontsize{20}{25}\selectfont\bfseries 本次账户资料及办理记录\par}
\vspace{7pt}{\noindent\small\color{Muted}仅供模拟使用·汽车贷款保险：已有保单与新增保障记录·13-03\par}\vspace{10pt}
\subsection*{A：提交已有保单并申请处理新增保障}
\textbf{已有个人车辆保单：} 继续保留


\textbf{已有保单的责任保障：} 包含对第三方人身或财产损害的责任保障


\textbf{已有保单的车辆损坏保障：} 包含碰撞损坏及综合损坏保障


\textbf{已有保单的贷款机构登记：} 已列明当前贷款机构的抵押权益


\textbf{新增保障的用途：} 对贷款机构的车辆抵押利益提供物理损坏保障


\textbf{新增保障的责任保障：} 不包含驾驶者对第三方的责任保障


\textbf{对新增保障的处理：} 上传正式保单证明，申请核对保障期间、取消对应新增保障并调整重叠期间费用


\textbf{需要准备的资料：} 保险声明页、有效期和车辆资料；保险公司可直接提供


\textbf{新增保障费用：} 核验通过后，按本次通知规定调整对应期间的费用


\textbf{后续保险登记：} 按银行要求提供续保或保单变更资料


\subsection*{B：保留新增保障及已有保单}
\textbf{已有个人车辆保单：} 同样继续保留


\textbf{已有保单的责任保障：} 包含对第三方人身或财产损害的责任保障


\textbf{已有保单的车辆损坏保障：} 包含碰撞损坏及综合损坏保障


\textbf{已有保单的贷款机构登记：} 已列明当前贷款机构的抵押权益


\textbf{新增保障的用途：} 对贷款机构的车辆抵押利益提供物理损坏保障


\textbf{新增保障的责任保障：} 不包含驾驶者对第三方的责任保障


\textbf{对新增保障的处理：} 不提交本次核对申请，新增保障及计费继续保留


\textbf{需要准备的资料：} 无本次上传资料操作


\textbf{新增保障费用：} 按现有安排计入贷款账户，并按贷款合同计息


\textbf{后续保险登记：} 原有保单与新增安排分别继续管理


\subsection*{办理确认事项}
\textbf{合同保险要求。} 银行要求的两种车辆损坏保障均在已有保单内；已有保单的保障金额和自付额条件已经核对符合贷款合同。贷款机构名称及车辆资料一致，当前也没有断保、欠保费或保单即将失效的情况。


\textbf{新增保障记录。} 本次新增安排是贷款机构在未登记有效保单证明时购入的抵押利益保障，不是你向原保险公司购买的升级项目。它不增加已有保单的责任保障，也不承担贷款差额保障、机械维修或延长保修服务。


\textbf{本次通知中的核对流程。} 银行接受正式保险声明页及承保证明，用来核对车辆、贷款机构和保障期间。证明覆盖新增保障的全部对应期间、且符合合同要求时，银行按通知流程取消该对应安排，并调整重叠期间收取的保障费及相关计息。本次无额外资料提交费。


\textbf{办理与付款。} 保险公司已经确认可以提供上述资料，你能够通过银行接受的渠道提交。办理期间继续保留个人保单，并按银行有效账单履行付款义务；申请本身不变更原贷款本金、利率或正常月供安排。


\textbf{本次决策范围。} 不更换个人保险公司，不取消已有保单，也不申请购买其他保险产品。选择的是是否用现有证明处理贷款账户上的新增保障记录。

\end{document}